\chapter{改进的规范作用量}

人们已经提出了多种方法，把规范作用量改进到 $O(a^2)$ 以上。其中包括 Symanzik 方法、平均场改进，以及受重整化群启发的完美作用量。让我以回顾 Symanzik 纲领来开始这一讨论。

\section{L\"uscher-Weisz 作用量}

所有威尔逊圈展开中的首阶项是
$\mathcal{O}^{(4)} = \sum_{\mu\nu} F_{\mu\nu}F_{\mu\nu}$，由于不存在维数 5 的算符，修正从 $O(a^2)$ 开始。因此，任何写成威尔逊圈线性组合的格点作用量都具有正确的连续极限，其修正为 $O(a^2)$ 阶。存在三个维数 6 的算符，用连续记号表示如下：
\begin{equation}
\begin{aligned}
\mathcal{O}_1^{(6)} &= \sum_{\mu,\nu} \Tr \bigg(D_\mu F_{\mu\nu}
    D_\mu F_{\mu\nu}\bigg),\\
\mathcal{O}_2^{(6)} &= \sum_{\mu,\nu,\rho} \Tr \bigg(D_\mu F_{\nu\rho}
    D_\mu F_{\nu\rho}\bigg),\\
\mathcal{O}_3^{(6)} &= \sum_{\mu,\nu,\rho} \Tr \bigg(D_\mu F_{\mu\rho}
    D_\nu F_{\nu\rho}\bigg).
\end{aligned}
\label{eq:gaugedimsix}
\end{equation}
在格点上也只能画出 3 个六链圈。如图~\ref{fig:7} 所示，它们分别是平面型的 $\mathcal{L}_1^{(6)}$、扭曲的 $\mathcal{L}_2^{(6)}$ 和 L 形的 $\mathcal{L}_3^{(6)}$。因此，格点作用量的经典改进，即消去 $O(a^2)$ 项，可以通过取方格和这三个六链圈的线性组合来实现。这些圈中的每一个都有展开式
\begin{equation}
\mathcal{L} = r^{(4)}\mathcal{O}^{(4)}
    + r_1^{(6)}\mathcal{O}_1^{(6)}
    + r_2^{(6)}\mathcal{O}_2^{(6)}
    + r_3^{(6)}\mathcal{O}_3^{(6)} + \ldots,
\label{eq:gaugeimpa}
\end{equation}

\begin{figure}[htbp]
\centering
\includegraphics[width=0.8\textwidth]{images/fig7.png}
\caption{六链圈的三种类型。}
\label{fig:7}
\end{figure}

L\"uscher 和 Weisz 给出了计算展开系数 $r_\alpha^{(d)}$ 的一个优雅方法 [73]。他们的结果汇总在表~\ref{tab:4} 中。

\begin{table}[htbp]
\centering
\begin{tabular}{lccccc}
\toprule
圈 & $r^{(4)}$ & $r_1^{(6)}$ & $r_2^{(6)}$ & $r_3^{(6)}$\\
\midrule
$\mathcal{L}^{(4)}$  & $-\frac{1}{4}$ & $\frac{1}{24}$ & 0 & 0\\
$\mathcal{L}_1^{(6)}$ & $-2$ & $\frac{5}{6}$ & 0 & 0\\
$\mathcal{L}_2^{(6)}$ & $-2$ & $-\frac{1}{6}$ & $\frac{1}{6}$ & $\frac{1}{6}$\\
$\mathcal{L}_3^{(6)}$ & $-4$ & $\frac{1}{6}$ & 0 & $\frac{1}{2}$\\
\bottomrule
\end{tabular}
\caption{4 链和 6 链威尔逊圈经典展开中维数 4 和维数 6 连续算符的系数。}
\label{tab:4}
\end{table}

格点规范作用量可以写为
\begin{equation}
\mathcal{S}_g = \frac{6}{g^2}\bigg\{c^{(4)}(g^2)\mathcal{L}^{(4)}
    + \sum_{i=1,3} c_i^{(6)}(g^2)\mathcal{L}_i^{(6)}\bigg\}
\label{eq:gaugeimpb}
\end{equation}
即用方格和这三个 6 链圈表示。利用表~\ref{tab:4} 中的结果，得到归一化条件
\begin{equation}
c^{(4)}(g^2) + 8c_1^{(6)}(g^2) + 8c_2^{(6)}(g^2)
    + 16c_3^{(6)}(g^2) = 1
\label{eq:gaugeimpc}
\end{equation}
使得在连续极限下作用量化为
$F_{\mu\nu}F_{\mu\nu}/4$。此外，从表~\ref{tab:4} 可以看出，树级改进可以通过如下选择获得
\begin{equation}
c^{(4)} + 20c_1^{(6)} = 0;\quad c_2^{(6)} = 0;\quad c_3^{(6)} = 0.
\label{eq:gaugeclassicalimp}
\end{equation}
这通过抵消两个平面圈的贡献消去 $\mathcal{O}_1^{(6)}$。在三个维数 6 的算符中，只有 $\mathcal{O}_1^{(6)}$ 破坏洛伦兹（转动）不变性。因此，重夸克 $q\bar q$ 势的转动不变性可以用来非微扰地调参 $c_1^{(6)}$，并比较各种改进作用量，这将在后文讨论。

为了纳入量子效应，Curci、Menotti 和 Paffuti 做了 1 圈计算。他们采用的标准是：所有威尔逊圈都在首阶对数水平得到改进 [74,75]。这对应于改进弦张力。（要在所有长度尺度上消去 $O(a^2)$ 项，而不仅仅是在长距离重夸克势中，就需要在 1 圈计算中包含有限项。完全改进的作用量，正如我们将在最终的 L\"uscher-Weisz 形式中看到的，将对 $c^{(4)}$ 和 $c_1^{(6)}$ 引入 $O(g^2)$ 修正，并需要 $\mathcal{O}_2^{(6)}$。）他们计算的结果要求改进作用量与式~\ref{eq:gaugeclassicalimp} 所给出的形式相同。与经典结果的这一致是偶然的。Symanzik 的在壳改进纲领一般性地规定，必须导出三个关系来固定 $c_i^{(6)}$，使得在壳量没有 $O(a^2)$ 修正。事实上，L\"uscher 和 Weisz 发现，改进所有大圈在树级只产生两个约束 [76]
\begin{equation}
\begin{aligned}
c^{(4)} + 20c_1^{(6)} - 4c_2^{(6)} + 4c_3^{(6)} &= 0\\
c_2^{(6)} &= 0.
\end{aligned}
\label{eq:gaugeimpd}
\end{equation}
因此，他们剩下一族 1-参数改进作用量。这一自由度的原因是存在由场重定义
$A_\mu \to A_\mu + a^2g^2d(g)D_\nu F_{\nu\mu}$ 产生的冗余算符 $\mathcal{O}_3^{(6)}$ [73,76]。由于在不影响在壳改进的前提下，可以选择属于冗余算符的某个相互作用的系数，方便的选择 $c_3^{(6)} = 0$ 简化了作用量，并再次把我们带回到式~\ref{eq:gaugeclassicalimp}。总之，式~\ref{eq:gaugeimpd} 是基于改进大威尔逊圈所能导出的最一般约束集。Curci 等人采用的标准，即改进 $q\bar q$ 势的长距离部分，是其子集。由于这涉及谱量，可以利用冗余算符得到更简单的形式式~\ref{eq:gaugeclassicalimp}。

关系 $c_3^{(6)} = 0$ 在树级和 1 圈级都成立，因为它来自对冗余算符的利用。将其与式~\ref{eq:gaugeimpd} 中的第一个约束以及归一化条件式~\ref{eq:gaugeimpc} 合并，得到一个条件，在树级为
\begin{equation}
c_1^{(6)} - c_2^{(6)} = -\frac{1}{12},
\end{equation}
这就是能从威尔逊圈得到的一切。为了得到第二个条件，L\"uscher 和 Weisz 不得不构造其他要改进的谱量。由于在真实理论中未能找到合适的条件，他们考察了一种胶子具有质量并表现为渐近态的情形。利用这些态的质量和散射振幅，他们导出了 1 圈改进作用量 [76]，这就是所谓的 L\"uscher-Weisz 作用量，其首阶修正为 $O(g^4a^4)$，
\begin{equation}
\begin{aligned}
c_0^{(4)}(g^2) &= \frac{5}{3} + 0.2370g^2\\
c_1^{(6)}(g^2) &= -\frac{1}{12} - 0.02521g^2\\
c_2^{(6)}(g^2) &= -0.00441g^2\\
c_3^{(6)}(g^2) &= 0.
\end{aligned}
\end{equation}
在这一 $c_i$ 形式中，归一化条件式~\ref{eq:gaugeimpc} 已被纳入。

关于 L\"uscher-Weisz 作用量的要点是：由于系数 $c_2^{(6)}$ 很小，该作用量在很好的近似下可以像经典情形那样，只保留 $\mathcal{L}^{(4)}$ 和 $\mathcal{L}_1^{(6)}$ 而得到改进。这个 1 圈结果可以毫不费力地进行平均场改进，如下文所述。

\section{蝌蚪图与平均场改进}

规范场用指数化形式
\begin{equation}
U_\mu(x) = e^{iagA_\mu(x)} = 1 + iagA_\mu(x)
    - \frac{a^2g^2}{2}A_\mu^2(x) + \ldots
\end{equation}
表示，会产生带有一个、两个、$\ldots$ 胶子的定域夸克-胶子顶点。在这些顶点中，除最低阶外都是格点赝象。收缩 $a^2g^2A_\mu^2(x)/2$ 中的两个胶子产生蝌蚪图。Parisi 以及 Lepage 和 Mackenzie [77,78] 认识到，这些赝象虽然按 $a$ 的幂次被朴素地压低，但实际上由于蝌蚪圈产生的紫外发散，只按 $g^2$ 的幂次被压低。他们论证，这些蝌蚪贡献是以下现象的根源：(i) 短距离量（如朗道规范中的链、小威尔逊圈、$\kappa_c$、$\ldots$）与其微扰估计之间的不匹配，以及 (ii) 微扰格点展开中的大系数。因此，他们提出了平均场改进（也称蝌蚪图改进(TI)）来消除这些赝象。

平均场改进首先假设格点场可以分成紫外(UV)部分和红外(IR)部分，并且紫外部分应该被积掉 [77,78]，即
\begin{equation}
e^{iagA_\mu(x)} = e^{iag(A_\mu^{IR}(x)+A_\mu^{UV}(x))}
    = u_0 e^{iagA_\mu^{IR}(x)} \equiv u_0 \tilde U_\mu(x).
\end{equation}
用一个整体常数因子 $u_0 < 1$ 重新标度链可以保持理论规范不变。在这一重新标度下（每个 $U \to u_0 \tilde U$），威尔逊规范作用和费米子作用量变为
\begin{equation}
\mathcal{S}^W \to \sum_x \bigg[\beta u_0^4\mathcal{W}_{11}
    + \bar{\tilde\psi}\tilde\psi
    - \kappa u_0 \bar{\tilde\psi}\tilde D\tilde\psi\bigg].
\end{equation}
因此，平均场改进的一个简化配方是使用更好的展开参数 $\tilde g^2 = g^2/u_0^4$、$\tilde\kappa =
\kappa u_0$、$\tilde U = U/u_0$ 和 $\tilde\psi =
\psi^{cont}/\sqrt{2\kappa u_0}$。注意，用带波浪号的变量表示时，理论看起来与原理论相同，只是所有链都被 $u_0$ 标度。Lepage 和 Mackenzie [78] 及其后的工作表明，对 1 圈估计做平均场改进在许多情况下给出的结果非常接近非微扰估计，而且展开系数要小得多。因此，蝌蚪图改进是一件值得做的事，尤其是因为在完成 1 圈计算之后只需极少的额外功夫。唯一的局限是它并非在所有情形下都同样有效，而且人们对它在某些情形下效果不佳的原因尚无真正的理解。为了说明蝌蚪图改进，我将其应用于 L\"uscher-Weisz 作用量。

蝌蚪因子 $u_0$ 有两种常见选择：(i) 方格的四次方根，和 (ii) 朗道规范中链的期望值。迄今为止大多数计算使用第一种选择，而 Lepage 及其合作者现在主张以朗道规范中的链作为更好的选择 [79]。要实现蝌蚪图改进，既需要这些变量的 1 圈表达式，也需要它们的非微扰值。对于威尔逊作用量
\begin{equation}
\begin{aligned}
u_0 &= \langle\mathcal{W}_{11}\rangle^{1/4} = 1 - \frac{1}{12}g^2\\
u_0 &= \langle U\rangle_{LG} = 1 - 0.0772g^2
\end{aligned}
\end{equation}
而对于 L\"uscher-Weisz 作用量，式~\ref{eq:gaugeclassicalimp}，则为 [80]
\begin{equation}
u_0 = \langle\mathcal{W}_{11}\rangle^{1/4}
    = 1 - \left(\frac{1}{2}\right)\left(\frac{4}{3}\right)
        \left(\frac{0.366262}{4}\right)g^2
    = 1 - 0.061044g^2.
\end{equation}

记 1 圈蝌蚪因子为 $u_0 = 1 - Xg^2$，则蝌蚪图改进的 L\"uscher-Weisz (TILW) 作用量为 [81]
\begin{equation}
\mathcal{S}_{TILW} = \beta_{eff}\bigg(\mathcal{L}^{(4)}
    - \frac{1}{20u_0^2}\big(1 + (0.1602-2X)g^2\big)\mathcal{L}_1^{(6)}
    - \frac{1}{u_0^2}0.00264g^2\mathcal{L}_2^{(6)}\bigg)
\label{eq:LWTIaction}
\end{equation}
$u_0$ 的非微扰值必须自洽地确定，因为作用量中的相对系数本身依赖于 $u_0$。注意，仅作为量子效应引入的 $\mathcal{L}_2^{(6)}$，在 TI 之后其系数被压低到可忽略的值。

关于 TILW 作用量值得注意的一个重要之点是 $\mathcal{L}_1^{(6)}$ 的强度。由于对当前使用的格点参数 $u_0 = 0.85-0.90$，$\mathcal{L}_1^{(6)}$ 的强度大于经典值 $1/20$，即蝌蚪图改进建议一种过修正。稍后我们将看到，这种过修正也是基于重整化群方法的改进作用量的特征。

对上述例子应用 TI 的一个简化描述是：$\mathcal{L}_1^{(6)}$ 的系数被乘上了 $u_0^2(pert)/u_0^2(non-pert)$，即方格与 $\mathcal{L}_1^{(6)}$ 之间多出的两条链的相对因子。在 $O(g^4)$ 处与 $u_0^2(non-pert)$ 不同的因子 $u_0^2(pert)$ 被吸收进 1 圈展开。显然，这使得 $g^2$ 的系数变小。因此，修正由可以极精确测量的 $u_0(non-pert)$ 主导。然而，由于选择 $u_0$ 的自由度，仍存在残余的不确定性。总之，TI 的基本思想是：大的蝌蚪贡献存在于所有量的微扰展开中，可以通过提取 $u_0(pert)$ 并按其链数取相应的幂次来吸收。这使得展开系数变小。然后，修改后的微扰展开用改进的 $\tilde g^2 = g^2/u_0^4$ 来表示，修正由整体因子 $u_0(non-pert)$ 主导，而后者应取为朗道规范中链的期望值。

用 TILW 作用量令人信服地展示改进的最早检验之一，是 $q\bar q$ 势中转动不变性的恢复。来自 [82] 的结果重绘在图~\ref{fig:8} 中。此后，利用有限温度相变 $T_c$ 的标度和拓扑磁化率的改进检验也已开展。正如稍后将讨论的，这些结果都强有力地支持改进。

\begin{figure}[htbp]
\centering
\includegraphics[width=0.8\textwidth]{images/fig8.png}
\caption{在格距 $a \approx 0.4$ fm 的 $6^4$ 格点上用威尔逊作用量和 TILW 作用量计算的静态夸克势。虚线是连续势的标准红外参数化 $V(r) = Kr - \pi/12r + c$，经调整以拟合势的轴上值。图重绘自 [82]。}
\label{fig:8}
\end{figure}

\section{重整化群与重整化轨迹}

这是对 Wilson 的重整化群表述及重整化轨迹概念的简要总结。细节见 [83--90]。我先从动量空间中的讨论开始。考虑在某一标度 $Q$ 处用广义作用量表述的有效理论
\begin{equation}
\mathcal{S}(Q) = \sum_\alpha K_\alpha^{(0)}(Q)S_\alpha
\end{equation}
其中 $S_\alpha$ 是所有可能的算符（相互作用），$K_\alpha^{(0)}$ 是相应的耦合常数。现在考虑一个精确重整化群变换 (RGT)
$\mathcal{T}_b$，它积掉 $Q$ 与 $Q/b \equiv
Q/2$ 之间的所有动量。这里 $b = 2$ 是变换的标度因子。标度 $Q/2$ 处的重整化理论为
\begin{equation}
\mathcal{S}(Q/2) = \sum_\alpha K_\alpha^{(1)}(Q/2)S_\alpha
\end{equation}
其中 $K_\alpha^{(1)}$ 是 $K_\alpha^{(0)}$ 的解析函数，
\begin{equation}
K_\alpha^{(1)} = \mathcal{T}_b\big[K_\alpha^{(0)}\big].
\end{equation}
反复使用 $\mathcal{T}_b$ 所定义的理论序列在耦合空间定义了一条流动轨迹。物理理论在每次 $\mathcal{T}_b$ 下都保持不变，从而沿整条轨迹保持不变。以格点单位测量的关联长度平凡地改变 $\xi/a \to \xi/ab$。该变换的不动点由下式定义
\begin{equation}
K_\alpha^{(*)} = \mathcal{T}_b\big[K_\alpha^{(*)}\big],
\end{equation}
即在 $\mathcal{T}_b$ 下理论再现自身的那些点。由于所有关联长度（动量）在 $\mathcal{T}_b$ 下都按 $b$ 标度，因此在不动点处
\begin{equation}
\xi\big|_{\mathrm{fixed\ point}} = 0 \quad \mathrm{or} \quad \infty.
\end{equation}
$\xi = 0$ 的不动点是平凡不动点。这类例子是大多数统计力学模型的 $T = 0$ 和 $T = \infty$ 极限。我们感兴趣的是 $\xi = \infty$ 的非平凡不动点。在这些点上，以格距为单位测量的关联长度 $\xi/a =
1/Ma$ 是无穷大。这使我们能够取格点场论的连续极限——固定 $M$ 不变令 $a = 0$。

现在考虑一个靠近不动点 $\{K_\beta^*\}$ 的点 $\{K_\beta\}$，以及一个 $b$ 无穷小的 RGT。在假设耦合常数空间中没有奇异性的前提下，$\mathcal{T}_b$ 产生的新耦合 $K_\alpha'(K_\beta)$ 可以在不动点附近做泰勒展开
\begin{equation}
K_\alpha'(K_\beta) = K_\alpha^* +
    \frac{\partial K_\alpha'}{\partial K_\beta}\bigg|_{K^*}\Delta K_\beta
    + \ldots.
\label{eq:ltma}
\end{equation}
在这种情况下，可以在 $\{K_\beta^*\}$ 的邻域内定义一个「线性区域」，其中
\begin{equation}
\Delta K_\alpha' =
    \frac{\partial K_\alpha'}{\partial K_\beta}\bigg|_{K^*}\Delta K_\beta
\label{eq:ltmb}
\end{equation}
这里 $\Delta K_\beta = K_\beta - K_\beta^*$。线性化变换矩阵（也称稳定性矩阵）
\begin{equation}
\mathcal{T}_{\alpha\beta} \equiv
    \frac{\partial K_\alpha'}{\partial K_\beta}\bigg|_{K^*}
\label{eq:ltmc}
\end{equation}
局部地控制着 RG 流动。在场论讨论中，稳定性矩阵 $\mathcal{T}_{\alpha\beta}$ 被称为
$\beta$-函数。注意，由于式~\ref{eq:ltma} 的截断，所有基于 $\mathcal{T}_{\alpha\beta}$ 分析的陈述都隐含着这样的假设：它们在靠近不动点的线性区域中且 $b$ 较小时近似成立，并且只有当 $\{K_\beta\} \to
\{K_\beta^*\}$ 时才严格成立。

Wilson 将实空间重整化群表述为上述动量空间重整化群的替代方案。这一方法的特点是对依赖于位置的原自由度 $s$ 做约束积分以产生有效理论，
\begin{equation}
e^{K_\alpha'S_\alpha(s')} = \int \mathcal{D}s\,\mathcal{T}(s',s)
    e^{K_\alpha S_\alpha(s)}.
\end{equation}
这里 $\mathcal{T}(s',s)$ 是规定如何对某个局域邻域（称为阻塞胞）中的 $s$ 取平均并用有效的 $s'$ 替换的约束。$\mathcal{T}(s',s)$ 上保持配分函数不变的必要条件是
\begin{equation}
\int \mathcal{D}s'\,\mathcal{T}(s',s) = 1.
\label{eq:blockcondition}
\end{equation}
可以看到，如果 $D$ 维中的阻塞胞为 $b^D$，并且约束 $\mathcal{T}(s',s)$ 等价于积分 $Q/2 - Q$ 范围内的动量，那么这两种方法是等价的。一般说来，对上半动量做精确积分是不可行的，$\mathcal{T}(s',s)$ 被构造成阻塞胞中原自旋的某种合理「平均」。除式~\ref{eq:blockcondition} 的条件外，合理「平均」的定义是它保持理论的物理性质。由于这种对偶描述（动量空间与实空间表述），RGT 也被称为阻塞变换——这一用法在统计力学和数值实现中更为常见。

下面对矩阵 $\mathcal{T}_{\alpha\beta}$ 的性质作一简要总结，该矩阵可以用动量空间或实空间表述来计算，其相关本征值 $\Lambda_i$ 和本征向量 $\Phi_i$ 由求解 $\mathcal{T}\Phi_i =
\Lambda_i\Phi_i$ 得到。

\begin{enumerate}
\item $\mathcal{T}_{\alpha\beta}$ 一般不是对称矩阵。然而其本征值是实的。

\item 本征值依赖于标度因子 $b$。由于两次相继变换 $b \otimes b$ 后得到的理论应等于一次按 $b^2$ 标度后得到的理论，本征值的一般形式为 $\Lambda_i = b^{\lambda_i}$，其中 $\lambda_i$ 与 $\mathcal{T}_b$ 无关，被称为临界指数。

\item 本征向量 $\Phi_i$ 按定义就是线性标度场。考虑沿某个本征向量 $\Phi$ 从不动点出发的大小为 $u$ 的小偏离。则在 RGT 下，耦合的重整化值 $u'$ 为
  \begin{equation}
  u'\Phi = \mathcal{T}_b u\Phi = b^{\lambda_\Phi}u\Phi.
  \label{eq:ltmd}
  \end{equation}
  注意，在推导 $u$ 的这一简单标度时关键性地用到了线性性假设。本征向量依赖于 $\mathcal{T}_b$ 的具体形式。

\item 本征值 $\Lambda_i > 1$ 导致沿这些 $\Phi_i$ 的偏离在 $\mathcal{T}_b$ 下增长，从而离开不动点。这些本征值及其对应的本征向量被称为相关的。

\item 对应本征值 $\Lambda_i < 1$ 的算符随阻塞步数按几何级数衰减。这些算符对系统的长距离性质没有贡献，因而被称为无关的。相关的指数 $\lambda_i$ 控制对标度的修正。

\item 恰为 1 的本征值称为边缘的。沿边缘方向的偏离在 RGT 下不变。要确认某个本征值真正边缘，必须超越线性近似。典型地，边缘算符在更高阶会发展出对数修正。快速考察渐近自由场论的
$\beta$-函数就能使你确信，$g$ 在树级是边缘耦合，并因量子修正而变得相关。

\item 还有一类本征向量，称为冗余算符，它们没有物理意义。冗余算符是作用量在重新定义场这一不变性下的体现，即它们定义不同算符之间的关系，例如运动方程。对应的本征值可以 $<1$、$=1$ 或 $>1$。这些本征值依赖于 RGT 的选择，即标度因子 $b$ 相同或不同的不同 RGT 会给出不同的 $\lambda_i$。由于不同的非线性 RGT 对应不同的场重新标度，冗余算符是变换特定的。事实上，把它们分离出来的一种方法是换一个不同的 RGT 重复计算 $\mathcal{T}_{\alpha\beta}$；随 RGT 改变的指数 $\lambda_i$ 对应冗余算符。尽管它们不携带任何物理信息，分离这些算符仍然很重要，尤其在数值计算中。例如，改进作用量时要利用所有可能的冗余算符来简化作用量，从而使计算成本最小化。

\item 在某些条件下 $\mathcal{T}_{\alpha\beta}$ 是块对角的，即耦合分解成在 RGT 下封闭的子集。例如，零磁场下 Ising 哈密顿量的自旋翻转对称性 $s \to -s$ 使相互作用分解成奇、偶两个扇区。因此，可以对每个扇区分别计算 $\mathcal{T}_{\alpha\beta}$，因为本征值谱是独立的。
\end{enumerate}

在这个广义耦合空间 $\{K_\alpha\}$ 中，存在称为临界点 $\{K_\alpha^c\}$ 的、具有 $\xi = \infty$ 的点。从临界点出发，RGT 产生另一个临界点，因为 $\xi' =
\xi/b = \infty$。因此，临界点集在无穷维空间 $\{K_\alpha\}$ 中定义了一个超曲面。该超曲面上的 RG 流动可以 (a) 随机漫游，(b) 趋向某个极限环或奇怪吸引子，或 (c) 收敛到某个不动点 $\{K_\alpha^*\}$。我们只关心可能性 (c)。在不动点处，重整化耦合精确等于原耦合 $K_\alpha' = K_\alpha$，理论在所有长度标度上再现自身。这与临界点不同，在临界点只有长距离行为被再现。不动点的位置依赖于 RGT，一般来说给定的 RGT 会有不止一个 $\xi = \infty$ 的不动点。每个这样的不动点都有一个吸引域，即在 RGT 下收敛到它的 $\{K_\alpha^c\}$ 集合。该超曲面因而正交于 $\mathcal{T}_b$ 的所有相关本征向量。吸引域定义了普适类，因为所有理论的长距离行为都相同，由不动点理论所支配。

沿以不动点为终点的相关标度场的流动称为重整化轨迹 (RT)。正如不动点是其普适类中所有临界点的吸引子，RT 是收敛于位于不动点吸引域内临界点的所有流动的吸引子。如图~\ref{fig:9} 所示。为简单起见，我假设只有一个相关耦合，正如 QCD 的情形。沿 RT 没有标度破坏，因为按构造所有无关标度场都为零。因此，用沿精确 RT 的作用量做模拟，将再现连续物理而无离散误差。RT 因而成为改进格点公式化的圣杯，即完美作用量。

\begin{figure}[htbp]
\centering
\includegraphics[width=0.8\textwidth]{images/fig9.png}
\caption{临界超曲面（$\xi = \infty$）、不动点、在阻塞变换下流向红外的重整化轨迹 (RT) 以及被它们吸引的流动的示意图。还显示了简单威尔逊作用量以及从它出发的流动。}
\label{fig:9}
\end{figure}

在有限格点上做模拟是在耦合空间中所有 $\xi \ll L$ 的区域进行的，而不是在不动点或临界点处。当前典型值处于 $\xi/a = 1/\sqrt{\sigma}a
= 5$--$15$ 的范围。即使对淬火模拟，这一范围在不久的将来也不会改变。该区域中的点位于终止于临界点的流动上，而这些临界点可能离不动点相当远。因此，在 $\xi \approx 10$ 处沿这些流动对标度的修正可能很大。可以通过调整作用量，使其位于另一条更靠近不动点出发、并在 $\xi \approx 10$ 处更接近 RT 的流动上来减小这些修正。因此，改进作用量的一种非微扰方法是估计 RT。由于流动被 RT 吸引，其位置可以通过从甚至简单的方格作用量出发、跟踪某种阻塞变换下产生的流动来估计。因此，沿流动计算阻塞作用量是所有改进作用量的重整化群方法的基础。确定阻塞作用量的方法见综述 [86]。

支撑重整化群有效性有一个标准假设，即不动点作用量是定域的（耦合的强度随威尔逊圈的大小以及被 Dirac 算符耦合的点之间的间隔按指数衰减），尽管如此，定域作用量仍然涉及「无穷」多个耦合。在评估阻塞产生的作用量以及做模拟时保留所有这些耦合是不切实际的。因此，人们不得不把作用量截断为有限（少量）个相互作用。此外，$\xi \approx 10$ 处的 RT 可能与不动点附近的 RT 在形式上差别很大。因此，设计改进作用量时要面对的问题是：(i) 保留的最小项集是什么，(ii) 是否所有冗余相互作用都已从该集中移除，以及 (iii) 如何比较用不同 RGT 评估的这类改进作用量。还要注意，RT 向截断子空间的投影可能并不是在该空间中能设计出的标度破坏最小的作用量。下面我将探讨回答这些问题的尝试。

我将通过回顾那些正在用于计算与唯象学相关量的改进规范作用量，来总结当前设计改进规范作用量的活动。不幸的是，不同的合作组没有使用标准的记号。为便于比较，我将尽量把它们写成统一的形式，同时保持与原始记号的联系。威尔逊圈用 $W^i_{\mu\nu}$ 表示，其中指标 $i$ 既用于圈的类型，也用于其表示。除非另有说明，隐含地取威尔逊圈迹的实部。
$W^i_{\mu\nu}$ 经过归一化，使迹的最大值为 1。如前所述，$\beta$ 是总耦合，$c_i$ 给出不同圈的相对强度。

\section{Iwasaki 作用量}

Iwasaki 在微扰论中研究了阻塞变换。他的标准是得到经过几次阻塞后接近 RT 的规范作用量。他的 Ans\"atz 为 [91]
\begin{equation}
\mathcal{S}_{Iwasaki} = \beta\sum_{\mu,\nu}\sum_x
    \bigg(W^{1\times1}_{\mu\nu} - 0.0907W^{1\times2}_{\mu\nu}\bigg)
\label{eq:iwasakiaction}
\end{equation}
其中 $W^{1\times2}$ 是图~\ref{fig:7} 所示的矩形圈 $\mathcal{L}_1^{(6)}$。（Iwasaki 使用了记号
$(1-8c_0)W^{1\times1}_{\mu\nu} + c_0W^{1\times2}_{\mu\nu}$，其中
$c_0 = -0.331$，这与式~\ref{eq:iwasakiaction} 在重新标度 $\beta$ 后相同。）该作用量表现出更好的转动不变性和 $T_c/\sqrt{\sigma}$ 的标度。JLQCD 合作组正在使用它 [92]。

\section{$b=\sqrt{3}$ 阻塞变换}

$\sqrt{3}$ 阻塞变换对 4 维规范理论具有一些非常好的性质，此外它的标度因子比常规的 $b = 2$ 更小 [89]。阻塞胞由一个格点及其八个近邻组成。因此，该变换对费米子和规范场都具有高得多的转动对称性 [89]。Gupta 和 Patel 用蒙特卡洛重整化群方法研究了阻塞下的流动。他们考察了 (i) 截断和 (ii) $c_i$ 随 $\beta$ 的变化 [87]。为了唯象学研究，他们把 RT 近似为四参数空间中的一条线性轨迹
\begin{equation}
\begin{aligned}
\mathcal{S}_{\sqrt{3}} = \beta\,\mathrm{Re}\Big\{&
    c_F^{1\times1}\sum\Tr W^{(1\times1)}
    + c_F^{1\times2}\sum\Tr W^{(1\times2)}\\
    &+ c_6^{1\times1}\sum\Big[\frac{3}{2}(\Tr W^{(1\times1)})^2
        - \frac{1}{2}\Tr W^{(1\times1)}\Big]\\
    &+ c_8^{1\times1}\sum\Big[\frac{9}{8}|\Tr W^{(1\times1)}|^2
        - \frac{1}{8}\Big]\Big\},
\end{aligned}
\label{eq:rootthreedef}
\end{equation}
其中
\begin{equation}
c_6^{1\times1}/c_F^{1\times1} = -0.12;\quad
c_8^{1\times1}/c_F^{1\times1} = -0.12;\quad
c_F^{1\times2}/c_F^{1\times1} = -0.04.
\label{eq:rootthreerat}
\end{equation}
这一特定的耦合选择（作用量的截断）基于以下原因。包含伴随表示中方格的耦合 $c_8^{1\times1}$ 是为了避免基础-伴随平面中的相结构（见第 15 节）；包含 6 维表示中的 $c_6^{1\times1}$ 则因为它只在 $O(a^4)$ 及更高阶才与 $c_8^{1\times1}$ 不同；包含基础表示中的 $1\times2$ 矩形 $c_F^{1\times2}$，则因为它是以 Symanzik 方式消去 $O(a^2)$ 项的最简单圈。

圈的经典展开，见式 7.3，给出归一化条件
\begin{equation}
\bigg[c_F^{1\times1} + \frac{5}{2}c_6^{1\times1}
    + \frac{9}{4}c_8^{1\times1}\bigg] + 8c_F^{1\times2} = 1,
\end{equation}
其中我把来自 $W^{(1\times1)}$ 的前三项合并为方格的有效贡献。利用式~\ref{eq:rootthreerat} 给出的比值，得到非常高的权重 $c^{(4)} = 1/0.11$，如 [51] 所示，这与抑制位错而对物理瞬子有利相关。类似地，利用表~\ref{tab:4}，首阶离散化修正 $\mathcal{O}_1^{(6)}$ 的系数为
\begin{equation}
\frac{1}{24}\bigg[c_F^{1\times1} + \frac{5}{2}c_6^{1\times1}
    + \frac{9}{4}c_8^{1\times1}\bigg]
    + \frac{5}{6}c_F^{1\times2}.
\end{equation}
对于式~\ref{eq:rootthreerat} 给出的比值，这不是零，而是被所提出的 $W^{(1\times2)}$ 强度过修正了。在 $1/a \approx 2$ GeV 处，这一过修正大于蝌蚪图改进的 L\"uscher-Weisz Ans\"atz，与
$\mathcal{S}_{Iwasaki}$ 相近，小于 $\mathcal{S}_{QCDTARO}$ 中的。

早期检验改进的谱研究和胶球谱测量 [93] 尚无定论。事后看来，未能量化改进是因为统计量不足，以及没有考察强耦合下转动不变性这类量——对简单威尔逊作用量，格点赝象在这些量上清晰可见。最近，这一作用量也被用于研究格点拓扑，结果表明短距离赝象被抑制了 [51]。

\section{$b=2$ 阻塞变换}

QCDTARO 合作组也研究了阻塞下的作用量流动 [94,95]。他们使用 $b = 2$ 阻塞变换，并用 Schwinger-Dyson 方程在阻塞格点上评估耦合。他们的总体结论是，阻塞作用量在很好的近似下可以用两个耦合 $W^{1\times1}$ 和 $W^{1\times2}$ 表示。他们对 RT 的估计不能用系数之间的线性关系来近似。在所模拟的耦合范围内，一个近似的 Ans\"atz 为
\begin{equation}
\mathcal{S}_{QCDTARO} = \beta\sum_i\sum_{\mu,\nu}\sum_x
    \bigg(W^{1\times1}_{\mu\nu} - 0.11W^{1\times2}_{\mu\nu}\bigg)
\end{equation}
注意，$W^{1\times2}$ 的系数大于蝌蚪图改进的 Luscher-Weisz、Iwasaki 或 $b = \sqrt{3}$ 作用量中的。

\section{经典完美作用量}

Hasenfratz 和 Niedermayer 证明，可以用在 $g = 0$ 附近对 RGT 做鞍点积分来导出渐近自由场论的经典完美作用量 [96,97]。他们使用 Swendsen 提出的 $b = 2$ 阻塞变换 [98]。这里的实际问题在于，鞍点积分给出的作用量是否可以用少数定域项来近似。将此方法应用于 SU(3) 以及对截断版本的调参仍在研究之中 [99]。当前版本的截断经典不动点作用量 (CFP) 为 [100]
\begin{equation}
\mathcal{S}_{CFP} = \beta\sum_{i,x}\{c_1^i(1 - \mathrm{Re}\,\Tr W^i)
    + \ldots + c_4^i(1 - \mathrm{Re}\,\Tr W^i)^4\}.
\label{eq:Perfectaction}
\end{equation}
其中对 $i$ 的求和涉及两类圈：方格，以及绕 3-立方体扭曲的 6 链圈（图~\ref{fig:7} 中的 $\mathcal{L}_2^{(6)}$）。圈的高次幂在鞍点积分中自然出现。由于强度迅速下降，作者选择截断到四次幂。选择这一特定的 6 链圈是为了使作用量在时间上保持定域，并有助于改进转动不变性。两类圈的系数 $c_1$--$c_4$ 的值见表~\ref{tab:5}。

\begin{table}[htbp]
\centering
\begin{tabular}{lccccc}
\toprule
圈 & $c_1$ & $c_2$ & $c_3$ & $c_4$\\
\midrule
$1\times1$        & 3.248   & $-1.580$ & 0.1257  & 0.0576\\
扭曲六链圈    & $-0.2810$ & 0.0051 & 0.0049  & $-0.0096$\\
\bottomrule
\end{tabular}
\caption{含 4 链和 6 链圈的截断「完美」作用量的当前系数。}
\label{tab:5}
\end{table}

目前尚未完全掌握关于截断稳定性的全面控制。第二个问题是，在 $a = 0$ 处得到的这一 Ans\"atz 在模拟所在的 $1/a \sim 2$ GeV 处是否仍是有用的参数化，即重整化轨迹是否近似线性。要检验这一点，需要在 $1/a \sim 2$ GeV 附近使用与鞍点积分中相同的 $b = 2$ 阻塞变换做非微扰调参。完成这一调参所需的计算尚未开始。因此，这一 Ans\"atz 被认为是进行中的工作。尽管如此，在 $T_c$ 标度、弦张力、拓扑磁化率以及转动对称性恢复方面的改进检验已经非常鼓舞人心 [100,52]。

\section{小结}

哪个是最改进的规范作用量这一问题还远未定论。显而易见，重整化群改进作用量的缺点是：展开是按圈的大小或表示的维数进行的，而不是按无关算符的维数进行的。由于每个圈对所有维数的算符都有贡献，如何参数化残余误差（$O(a^2) + O(a^4)$，或 $O(a^4)$，或更高）并不清楚。此外，实现运动方程（消除与冗余方向相关的自由度）也并非易事。因此，对于应该保留的明显最小算符集尚无共识。人们基于以下考虑添加了不同的算符：(i) 它们在微扰展开中的重要性，(ii) 阻塞作用量中系数的大小（记住，选择阻塞变换的自由度以及由此产生的相关 RT 会影响这些系数），以及 (iii) 具体纳入非微扰效应；例如，包含一个负的伴随耦合有助于避免基础-伴随平面中的相结构（见第 15 节），添加模拟 $F^{\mu\nu}\tilde F_{\mu\nu}$ 算符的圈有助于给出标度不变的瞬子 [100,101]。人们希望未来能出现更精确的调参规范作用量的方法。幸运的是，尽管存在这种模糊性，形势非常令人鼓舞。在转动不变性的恢复（用 $q\bar
q$ 势测量，如图~\ref{fig:8} 所示）、胶球谱以及有限温度研究中所看到的改进非常显著。在很好的近似下，这些作用量中的大多数都可以很好地用仅添加 $W^{(1\times2)}$ 项（其系数大于蝌蚪图改进所建议的值）并保留方格的少数几个高表示来表征。这样的截断空间使计算问题保持可处理。

要点是：尽管没有普遍接受的改进规范作用量，上述任何一种作用量所看到的改进都是显著的。目前，就误差大小和计算成本而言，更紧迫的挑战是改进 Dirac 离散化。这将在下一节讨论。
